% Preambolo comune dell'allegato tecnico. Si compila con XeLaTeX + BibTeX.
\documentclass[10pt,a4paper]{article}
\usepackage[a4paper,margin=2.2cm]{geometry}
\usepackage{fontspec}
\setmainfont{Fira Sans}[Scale=0.95,BoldFont={Fira Sans SemiBold},ItalicFont={Fira Sans Italic}]
\setmonofont{DejaVu Sans Mono}[Scale=0.82]
\usepackage{polyglossia}\setmainlanguage{italian}
\usepackage{amsmath,amssymb,amsthm,mathtools}
\usepackage{xcolor,tikz,booktabs,tabularx,array,enumitem,microtype,ragged2e,caption,float,longtable}
\usetikzlibrary{positioning,arrows.meta,calc,angles,quotes,patterns,decorations.pathreplacing}
\usepackage[numbers,sort&compress]{natbib}
\usepackage[hidelinks]{hyperref}
\definecolor{accent}{HTML}{0B5C8A}
\definecolor{muted}{HTML}{55606C}
\definecolor{boxbg}{HTML}{EEF3F7}
\setlength{\parindent}{0pt}\setlength{\parskip}{4pt}
\setlist{leftmargin=1.4em,itemsep=1.5pt,topsep=2pt}
\captionsetup{font=small,labelfont=bf}
\newcolumntype{L}{>{\RaggedRight\arraybackslash}X}
% Ambienti numerati per capitolo (sezione)
\theoremstyle{plain}
\newtheorem{proposizione}{Proposizione}[section]
\newtheorem{lemma}[proposizione]{Lemma}
\newtheorem{corollario}[proposizione]{Corollario}
\theoremstyle{definition}
\newtheorem{definizione}[proposizione]{Definizione}
\newtheorem{ipotesi}[proposizione]{Ipotesi}
\newtheorem{controesempio}[proposizione]{Controesempio}
\theoremstyle{remark}
\newtheorem{osservazione}[proposizione]{Osservazione}
% Livello di prova di un risultato: \prova{D} dimostrato, {C} calcolo, {S} simulazione, {P} dati pubblici, {N} non verificato sperimentalmente
\newcommand{\prova}[1]{\textsf{\footnotesize\color{accent}[#1]}}
% Riquadro per condizioni di validità e limiti
\usepackage{framed}\colorlet{shadecolor}{boxbg}
\newenvironment{riquadro}[1]{\begin{shaded}\textbf{#1}\par\smallskip}{\end{shaded}}
% Notazione comune
\newcommand{\Reali}{\mathbb{R}}
\newcommand{\vmax}{v_{\max}}
\newcommand{\tocc}{\tau_{\mathrm{occ}}}
\newcommand{\Oss}{\mathcal{O}}      % insieme osservato libero
\newcommand{\Omb}{\Omega}           % ombra
\newcommand{\Pres}{\Pi}             % possibile presenza
\newcommand{\Reg}{\mathcal{R}}      % regione di previsione
\newcommand{\unita}[1]{\,\mathrm{#1}}
\DeclareMathOperator{\dilata}{dil}
